%root = ../thesis-main.tex
\chapter{Implementation and Evaluation}\label{ch:impl_eval}

\section{The Reactive Framework}
Building upon the design principles established in
\Cref{ssec:reactive_alchemist_principles} and the architectural analysis in
\Cref{ssec:in_house}, this section presents the concrete implementation of the
proposed reactive solution. The framework is designed around a lightweight,
functional observer pattern that prioritises composition and granular control
over state changes.

\subsection{An Enhanced Observer Pattern}
The core abstraction of the framework is the \kotlin{Observable<T>} interface.
Unlike standard observer implementations that often enforce rigid class
hierarchies, this framework defines a minimal contract focused on the
subscription mechanism.

\code{Kotlin}{listings/Observable.kt}{The simplified Observable contract}{observable}

As shown in \Cref{lst:observable}, the \kotlin{onChange} method requires a
\emph{registrant} (of type \kotlin{Any}) alongside the callback function. This
design choice decouples the identity of the subscriber (herein referred to
interchangeably as \emph{subscribers}, \emph{registrants}, \emph{observers} or
\emph{listeners}) from the behaviour defined in the callback, allowing any
component to manage its subscriptions explicitly without implementing specific
listener interfaces. While the interface extends \kotlin{Disposable} to support
resource management, the implications of lifecycle management and the
prevention of memory leaks will be detailed in \Cref{ssec:disposability}.

To enforce strict control over state changes, the framework applies the
\emph{Interface Segregation Principle} (ISP) by splitting read and write
capabilities. The base \kotlin{Observable<T>} provides a read-only view of the
data stream, ensuring that consumers cannot inadvertently alter the state of
the source. State mutation is reserved for the \kotlin{MutableObservable<T>}
subtype. This interface extends the contract by exposing an \kotlin{update}
function, which encourages functional state transitions of the form \kotlin{(T)
-> T}. This distinction allows the internal logic of a component to mutate
state safely while exposing only the immutable \kotlin{Observable} interface to
external observers, preserving encapsulation.

\paragraph{Monadic Composition and Data Flow}
A primary advantage of this reactive approach is the ability to model complex
behaviours through the composition of simple event sources. Drawing inspiration
from functional programming and category theory, the framework treats
\kotlin{Observable} as a \emph{Monad}.

Through the basic uplifting operation \kotlin{fun observe(value: T):
Observable<T>} and the set of extension functions defined in the framework,
developers can chain operations using standard functional operators:
\begin{itemize}
	\item \texttt{map} which transforms the emitted values of an observable,
		creating a new derived stream;
	\item \texttt{flatMap}/\texttt{switchMap} with the goal to project each
		emitted value into a new observable and flatten the resulting streams,
		allowing for dynamic dependency replacement.
\end{itemize}
This monadic structure allows for the definition of complex data flows where
dependencies are handled implicitly. For example, a reaction's firing rate can
be defined as a transformation of local concentration values: as the
concentrations change, the derived rate updates automatically, propagating the
new value down the chain.

\subsection{Reactive Collections and Aggregation}
While standard observables handle scalar values efficiently, simulating a
pervasive ecosystem requires tracking aggregates, such as the list of
neighbours or the map of molecules within a node. To address this, the
framework introduces specific reactive collections: \kotlin{ObservableList},
\kotlin{ObservableSet}, and \kotlin{ObservableMap}.
These interfaces provide granular observation capabilities. For instance,
\kotlin{ObservableMap} allows observing the value of a specific key, while
\kotlin{ObservableSet} allows observing the membership status of a specific
element. However, the true power of these collections lies in their ability to
perform \emph{complex aggregation} over dynamic content.

\paragraph{Aggregating Dynamic Streams}
In a simulation, an aggregate value (e.g., the total concentration of a
molecule across all neighbours) depends on two distinct factors:
\begin{enumerate*}[label=(\roman*)]
	\item \emph{structural changes} with neighbors entering or leaving the list
		and 
	\item \emph{element changes} where the concentration value changing within
		an existing neighbor.
\end{enumerate*}

To handle this dual dependency, the framework provides sophisticated operators
such as \kotlin{combineLatest} for lists. This operator transforms an
\kotlin{ObservableList<T>} into a single \kotlin{Observable<R>} by applying a
mapping function \kotlin{(T) -> Observable<S>} to each element and then
aggregating the results. Crucially, the resulting stream emits a new value
whenever the list structure changes \emph{or} whenever any of the inner
observables emit a change. This allows to express through Alchemist global
properties, such as field strength or crowd density as continuous signals
derived purely from the local state of individual nodes.

The names and behaviours of such operators is largely inspired by \textsc{ReactiveX}
operators\footnote{\url{https://reactivex.io/documentation/operators.html}}.

\subsection{Handling Optionality: Native Null-Safety versus Monadic Options}
A notable detail in the implementation of the reactive framework is the
adoption of the \texttt{arrow.core.Option} type over Kotlin's built-in nullable
types (\texttt{T?}).

In a pure Kotlin environment, the language's native null-safety features (via
the \texttt{?} operator and smart casts) are typically preferred for their
zero-overhead runtime representation and syntactic convenience. However, the
Alchemist simulator operates as a polyglot environment, integrating components
written in Java, Kotlin and Scala (e.g., the ScaFi incarnation). This
necessitates a more robust strategy for handling value absence for two primary
reasons:

\begin{enumerate}
	\item \emph{Semantic Distinction in Streams:} In a reactive data flow, it
		is often necessary to distinguish between ``the value is known to be
		null'' and ``there is no value yet'' (or ``the value has been
		removed''). Kotlin's \texttt{T?} collapses these states. For instance,
		in \texttt{ObservableMap}, retrieving the value for a specific key
		returns an \texttt{Observable<Option<V>>}. This signature allows the
		stream to explicitly signal three distinct states:
    \begin{itemize}
		\item \kotlin{Some(v)}: The key is present with value $v$;
		\item \kotlin{Some(null)}: The key is present, and the value is
			explicitly null (if $V$ is nullable);
		\item \kotlin{None}: The key is missing from the map.
    \end{itemize}

	\item \emph{JVM Interoperability:} While Kotlin's nullable types are
		effective within Kotlin sources, they are represented as standard
		references (with metadata annotations) at the bytecode level. When
		accessed from Java or Scala, the strict compile-time safety guarantees
		are often lost or require verbose defensive programming. By using an
		explicit monadic wrapper like \texttt{arrow.core.Option}, we expose a
		consistent \ac{API} across all \ac{JVM} languages. This aligns with the
		functional paradigms of Scala (which relies heavily on its own
		\texttt{Option}) and provides a richer, composable interface (e.g.,
		\texttt{map}, \texttt{flatMap}, \texttt{fold}) than standard Java null
		checks.
\end{enumerate}

Consequently, while the use of an external library introduces a dependency, the
semantic precision and interoperability benefits of \texttt{arrow.core.Option}
outweigh the simplicity of native nulls in this specific architectural context.

\subsection{Optimising derived-observable values}
While the \kotlin{Observable<T>} interface defines the contract for data
consumption, the framework's efficiency lies in how these streams are
transformed and combined. The implementation of the aforementioned operators
like \kotlin{map}, \kotlin{filter} and \kotlin{mergeWith} relies on the
abstraction of \kotlin{DerivedObservable<T>}. This class manages the complexity
of state synchronisation and lifecycle management specifically addressing the
tension between the \emph{eager} consistency required by the \ac{DES} scheduler
and the desire for \emph{lazy} computation to minimise overhead and improve
resources usage.

\paragraph{The Challenge of Simulation Initialisation}
Alchemist divides its lifecycle into two distinct phases. The first is the
\emph{set-up phase}, where the environment is instantiated, nodes are created,
and reactions are wired to conditions. The second is the \emph{execution
phase}, where the simulation loop begins.

This separation poses a challenge for reactive programming. When defining a
reaction during the set-up phase, a programmer might attempt to create an
observation stream on a component that is not yet fully initialised. For
instance, observing a node's position before the node has been deployed into
the environment will trigger a runtime error if the observation is evaluated
eagerly.

To address this, the framework offers two strategies:
\begin{enumerate}
    \item \emph{Deferred Creation:} Creating eager event streams only within
      the \kotlin{initializationComplete} method of the \kotlin{Actionable}
      interface. This guarantees that the stream is constructed only after the
      set-up phase concludes and all components are stable;
    \item \emph{Lazy Subscription:} Creating event streams in the constructor
      but deferring their activation. The stream remains inert during the
      set-up phase and begins emitting values only when the simulation starts.
\end{enumerate}
To support the latter approach, the \kotlin{onChange} signature includes an
\kotlin{invokeOnRegistration} flag. When set to \kotlin{false}, it signals that
the registration should not trigger an immediate propagation of the current
value. This allows the dependency graph to be constructed safely during the
set-up phase without triggering premature computations on uninitialised data.

\paragraph{Smart Activation in Derived Observables}
The \kotlin{DerivedObservable<T>} class leverages this mechanism to implement a
``smart'' activation strategy that optimises both the computational graph and
resource usage.
A naive implementation of a reactive operator (e.g., \kotlin{map}) might simply
forward every downstream subscription to the upstream source. In a system with
many observers, this would lead to a \emph{subscription explosion}, where the
upstream source manages an excessive number of redundant listeners.

Instead, \kotlin{DerivedObservable} acts as a \emph{multiplexer}. It maintains
a single active subscription to its upstream source, regardless of how many
downstream observers are registered. When the upstream emits a value, the
derived observable computes the transformation once and multicasts the result
to all its subscribers. This significantly reduces the cardinality of the
dependency graph and the overhead of listener traversal.

Crucially, \kotlin{DerivedObservable} is \emph{inert} by default. It does not
subscribe to its upstream source until at least one observer registers with it.
If a derived stream is defined but never observed, it consumes no CPU cycles
and attaches no listeners to the source. This behaviour is managed via the
\kotlin{startMonitoring} and \kotlin{stopMonitoring} lifecycle hooks, which are
automatically invoked based on the emptiness of the subscribers list.

\paragraph{Explicit Initialization and Safety}
Finally, the framework addresses the semantic ambiguity of "initial values" in
push-based systems. As highlighted in \cite{bainomugisha2013SRP}:
\begin{quote}
    ``In push-based approaches, the programmer must initialise behaviours
    explicitly to make sure they hold a value when eagerly evaluating code in
    which they are used.''
\end{quote}

To strictly adhere to this principle, \kotlin{DerivedObservable} wraps its
internal state in an \kotlin{Option<T>}. This allows the framework to
distinguish between ``the stream has not yet yielded a value'' (\kotlin{arrow.core.None}) and ``the
stream yielded a null value'' (\kotlin{Some(null)}).

This distinction is vital for complex aggregators like \kotlin{combineLatest}.
If a derived stream attempts to aggregate an empty collection of observables,
it cannot compute a valid fresh value. Instead of blocking or throwing an
exception, the framework returns \kotlin{None}. This signals to downstream
consumers that the behaviour is defined but currently holds no value, allowing
the simulation to proceed without glitches until the necessary data becomes
available.


\subsection{\texttt{Disposable} pattern and Resource Management through Lifecycles}\label{ssec:disposability}
